% answers.tex —— 1954 年·全国卷 · 参考答案与详细解析片段
% 由 tools/generate_projects.py 依来源总账自 1954/1954_数学真题_试题.md + 1954/1954_数学真题_解析正文.md + 1954/1954_数学真题_答案要点.md 生成，勿手改（改源 MD 后重跑）
\gkpart{一、小题}

\begin{answer}
\gksolhead{1}{}

\textbf{题目}：

（甲）化简 $\left[\left(a^{\frac32}\cdot b^{-2}\right)\cdot\left(a\cdot b^{-3}\right)^{\frac12}\cdot\left(b^{\frac12}\right)^7\right]^{\frac13}$．

\Needspace{4\baselineskip}
\textbf{解析}

\Needspace{5\baselineskip}
逐因子展开并合并同底数幂：

\nopagebreak[4]
\[
\left(a^{\frac32}\cdot b^{-2}\right)\cdot\left(a\cdot b^{-3}\right)^{\frac12}\cdot\left(b^{\frac12}\right)^7
=a^{\frac32}\cdot b^{-2}\cdot a^{\frac12}\cdot b^{-\frac32}\cdot b^{\frac72}
=a^{\frac32+\frac12}\cdot b^{-2-\frac32+\frac72}.
\]

\Needspace{5\baselineskip}
$a$ 的指数为 $2$，$b$ 的指数为 $0$，故括号内为 $a^2b^0=a^2$，再开三次方得

\nopagebreak[4]
\[
\left(a^2\right)^{\frac13}=a^{\frac23}.
\]

\gkanswer{答案：$a^{\frac23}$}

\end{answer}
\begin{answer}
\gksolhead{2}{}

\textbf{题目}：

（乙）解 $\dfrac16\log x=\dfrac13\log a+2\log b+\log c$．

\Needspace{4\baselineskip}
\textbf{解析}

\Needspace{5\baselineskip}
两边同乘 $6$：

\nopagebreak[4]
\[
\log x=2\log a+12\log b+6\log c=\log\left(a^2b^{12}c^6\right).
\]

故 $x=a^2b^{12}c^6$。

\gkanswer{答案：$x=a^2b^{12}c^6$}

\end{answer}
\begin{answer}
\gksolhead{3}{}

\textbf{题目}：

（丙）用二项式定理计算 $(3.02)^4$，使误差小于 $\dfrac1{1000}$．

\Needspace{4\baselineskip}
\textbf{解析}

\Needspace{5\baselineskip}
把 $3.02=3+0.02$ 代入二项展开：

\nopagebreak[4]
\[
(3.02)^4=(3+0.02)^4=3^4+4\cdot3^3\cdot0.02+6\cdot3^2\cdot0.02^2+4\cdot3\cdot0.02^3+0.02^4.
\]

\Needspace{5\baselineskip}
第 $4$ 项 $4\cdot3\cdot0.02^3=0.000096<0.001$，故从第 $4$ 项起可舍去，误差小于 $\dfrac1{1000}$。前两项与第三项：

\nopagebreak[4]
\[
(3.02)^4\approx81+4\cdot27\cdot0.02+6\cdot9\cdot0.0004=81+2.16+0.0216=83.1816\approx83.182.
\]

\gkanswer{答案：$83.182$}

\end{answer}
\begin{answer}
\gksolhead{4}{}

\textbf{题目}：

（丁）直角三角形弦上半圆的面积等于勾上半圆与股上半圆面积之和，试证明之．

\Needspace{4\baselineskip}
\textbf{解析}

\Needspace{5\baselineskip}
设直角三角形的勾、股、弦分别为 $a,b,c$，则 $c^2=a^2+b^2$。弦上半圆的面积为

\nopagebreak[4]
\[
\frac12\pi\left(\frac c2\right)^2=\frac{\pi c^2}{8}.
\]

\Needspace{5\baselineskip}
勾、股上半圆的面积之和为

\nopagebreak[4]
\[
\frac12\pi\left(\frac a2\right)^2+\frac12\pi\left(\frac b2\right)^2
=\frac{\pi(a^2+b^2)}{8}.
\]

由 $c^2=a^2+b^2$ 知两者相等，即弦上半圆的面积等于勾上半圆与股上半圆面积之和。

\gkanswer{答案：见解析版}

\end{answer}
\begin{answer}
\gksolhead{5}{}

\textbf{题目}：

（戊）已知球的半径为 $r$，求内接正方体的体积．

\Needspace{4\baselineskip}
\textbf{解析}

\Needspace{5\baselineskip}
内接正方体的中心即球心，过中心的对角线为球的直径，长为 $2r$。设正方体边长为 $a$，其空间对角线为 $\sqrt3\,a$，故

\nopagebreak[4]
\[
\sqrt3\,a=2r\implies a=\frac{2r}{\sqrt3}=\frac{2\sqrt3}{3}r.
\]

\Needspace{5\baselineskip}
体积

\nopagebreak[4]
\[
V=a^3=\left(\frac{2\sqrt3}{3}r\right)^3=\frac{8\cdot3\sqrt3}{27}r^3=\frac{8\sqrt3}{9}r^3.
\]

\gkanswer{答案：$\dfrac{8\sqrt3}{9}r^3$}

\end{answer}
\begin{answer}
\gksolhead{6}{}

\textbf{题目}：

（己）已知三角形的一边之长为 $a$，两邻角为 $\beta$ 及 $\gamma$，求计算边长 $b$ 的计算公式．

\Needspace{4\baselineskip}
\textbf{解析}

\Needspace{5\baselineskip}
边 $a$ 所对的角为 $180^\circ-(\beta+\gamma)$。由正弦定理

\nopagebreak[4]
\[
\frac a{\sin(\beta+\gamma)}=\frac b{\sin\beta}
\implies b=\frac{a\sin\beta}{\sin(\beta+\gamma)}.
\]

\gkanswer{答案：$b=\dfrac{a\sin\beta}{\sin(\beta+\gamma)}$}

\par\medskip

\end{answer}
\gkpart{二、大题}

\begin{answer}
\gksolhead{7}{}

\textbf{题目}：

（二）描绘 $y=3x^2-7x-1$ 之图象，并按下列条件分别求 $x$ 的值的范围：（i）$y>0$；（ii）$y<0$．

\Needspace{4\baselineskip}
\textbf{解析}

\Needspace{5\baselineskip}
配方：

\nopagebreak[4]
\[
y=3x^2-7x-1=3\left(x-\frac76\right)^2-\frac{61}{12}.
\]

\Needspace{5\baselineskip}
故顶点为 $\left(\dfrac76,-\dfrac{61}{12}\right)$，抛物线开口向上。与 $x$ 轴的交点满足 $3x^2-7x-1=0$，解得

\nopagebreak[4]
\[
x=\frac{7\pm\sqrt{49+12}}{6}=\frac{7\pm\sqrt{61}}{6}.
\]

因此：

\begin{itemize}
\item $y>0$（图象在 $x$ 轴上方）时，$x<\dfrac{7-\sqrt{61}}6$ 或 $x>\dfrac{7+\sqrt{61}}6$；
\item $y<0$（图象在 $x$ 轴下方）时，$\dfrac{7-\sqrt{61}}6<x<\dfrac{7+\sqrt{61}}6$。
\end{itemize}

\gkanswer{答案：$y>0$ 时 $x\in\left(-\infty,\dfrac{7-\sqrt{61}}6\right)\cup\left(\dfrac{7+\sqrt{61}}6,+\infty\right)$；$y<0$ 时 $x\in\left(\dfrac{7-\sqrt{61}}6,\dfrac{7+\sqrt{61}}6\right)$}

\end{answer}
\begin{answer}
\gksolhead{8}{}

\textbf{题目}：

（三）假设两圆互相外切，求证用连心线段为直径所作的圆必与前两圆的外公切线相切．

\Needspace{4\baselineskip}
\textbf{解析}

设 $\odot O_1$ 与 $\odot O_2$ 外切，一条外公切线切两圆于 $A_1,A_2$。取连心线 $O_1O_2$ 的中点 $O$，过 $O$ 作 $OA\perp A_1A_2$。

\Needspace{5\baselineskip}
由于 $O_1A_1\perp A_1A_2$、$O_2A_2\perp A_1A_2$，且 $O_1A_1\parallel O_2A_2\parallel OA$，在梯形 $O_1A_1A_2O_2$ 中 $O$ 是 $O_1O_2$ 的中点，故 $OA$ 是梯形中位线，

\nopagebreak[4]
\[
OA=\frac12(O_1A_1+O_2A_2)=\frac12(r_1+r_2).
\]

又两圆外切，圆心距 $O_1O_2=r_1+r_2$，故 $OA=\dfrac12 O_1O_2$。于是以 $O$ 为圆心、$\dfrac12 O_1O_2$ 为半径（即以 $O_1O_2$ 为直径）的圆，其圆心到 $A_1A_2$ 的距离恰等于半径，故与 $A_1A_2$ 相切。另一条外公切线同理可证。

\gkanswer{答案：见解析版}

\end{answer}
\begin{answer}
\gksolhead{9}{}

\textbf{题目}：

（四）解 $\dfrac{1+\tan x}{1-\tan x}=1+\sin2x$，求 $x$ 之通值．

\Needspace{4\baselineskip}
\textbf{解析}

\nopagebreak[4]
\[
\frac{1+\tan x}{1-\tan x}=\frac{\cos x+\sin x}{\cos x-\sin x},
\]

\Needspace{5\baselineskip}
而 $1+\sin2x=\sin^2x+\cos^2x+2\sin x\cos x=(\sin x+\cos x)^2$。故原方程化为

\nopagebreak[4]
\[
\frac{\cos x+\sin x}{\cos x-\sin x}=(\cos x+\sin x)^2
\implies (\cos x+\sin x)\left[1-(\cos x-\sin x)(\cos x+\sin x)\right]=0.
\]

\nopagebreak[4]
\[
(\cos x-\sin x)(\cos x+\sin x)=\cos^2x-\sin^2x=\cos2x,
\]

\Needspace{5\baselineskip}
故

\nopagebreak[4]
\[
(\cos x+\sin x)(1-\cos2x)=0
\implies 2(\cos x+\sin x)\sin^2x=0.
\]

由 $\cos x+\sin x=0$ 得 $\tan x=-1$，$x=k\pi-\dfrac\pi4$；由 $\sin x=0$ 得 $x=k\pi$。经检验两者均为通解。

\gkanswer{答案：$x=k\pi-\dfrac\pi4$ 或 $x=k\pi$（$k$ 为整数）}

\end{answer}
\begin{answer}
\gksolhead{10}{}

\textbf{题目}：

（五）有一直圆锥，全面积为 $a$；与之同底同高之直圆柱全面积为 $a'$．求该圆锥高与母线之比．

\Needspace{4\baselineskip}
\textbf{解析}

\Needspace{5\baselineskip}
设圆锥的高为 $h$、底面半径为 $R$、母线为 $l$，则 $l^2=h^2+R^2$。圆锥全面积

\nopagebreak[4]
\[
a=\pi R(R+l),
\]

\Needspace{5\baselineskip}
同底同高的圆柱全面积

\nopagebreak[4]
\[
a'=2\pi R(R+h).
\]

\Needspace{5\baselineskip}
于是

\nopagebreak[4]
\[
\frac a{a'}=\frac{R+l}{2(R+h)}
\implies 2a(R+h)=a'(R+l).
\]

\Needspace{5\baselineskip}
代入 $R=\sqrt{l^2-h^2}$：

\nopagebreak[4]
\[
2a\left(\sqrt{l^2-h^2}+h\right)=a'\left(\sqrt{l^2-h^2}+l\right),
\]

\nopagebreak[4]
\[
(2a-a')\sqrt{l^2-h^2}=a'l-2ah.
\]

\Needspace{5\baselineskip}
两边同除以 $l$ 并令 $t=\dfrac hl$：

\nopagebreak[4]
\[
(2a-a')\sqrt{1-t^2}=a'-2at.
\]

\Needspace{5\baselineskip}
两边平方：

\nopagebreak[4]
\[
(2a-a')^2(1-t^2)=a'^2-4aa't+4a^2t^2,
\]

\Needspace{5\baselineskip}
整理得关于 $t$ 的一元二次方程

\nopagebreak[4]
\[
\left(8a^2-4aa'+a'^2\right)t^2-4aa't+\left(4aa'-4a^2\right)=0.
\]

\Needspace{5\baselineskip}
其判别式

\nopagebreak[4]
\[
\Delta=(-4aa')^2-4(8a^2-4aa'+a'^2)(4aa'-4a^2)=16a(2a-a')^3>0.
\]

\Needspace{5\baselineskip}
解得

\nopagebreak[4]
\[
t=\frac hl=\frac{2aa'\pm2(2a-a')\sqrt{a(2a-a')}}{4a^2+(2a-a')^2}.
\]

\gkanswer{答案：$\dfrac hl=\dfrac{2aa'\pm2(2a-a')\sqrt{a(2a-a')}}{4a^2+(2a-a')^2}$}

\end{answer}